\documentclass[10pt]{article}
\usepackage[utf8]{inputenc}
\usepackage[T1]{fontenc}
\usepackage[spanish]{babel}
\usepackage[letterpaper,margin=1.8cm]{geometry}
\usepackage{longtable,booktabs,array}
\usepackage[hidelinks]{hyperref}
\setlength{\parindent}{0pt}
\setlength{\parskip}{6pt}
\begin{document}
\begin{center}\Large\textbf{Respuesta actualizada a las observaciones del tutor}\end{center}
\textbf{Corte de avances: 15 de octubre de 2026.}

La versión actualizada se basa en los archivos experimentales disponibles. El corpus ampliado está construido con 96 pares; las ejecuciones analizadas usan 42 pares. La ampliación del corpus no se presenta como una ejecución ya realizada. Los originales se conservan y las referencias siguientes corresponden a las secciones de la versión de avances.

\small
\begin{longtable}{@{}p{0.5cm}p{3.2cm}p{2.5cm}p{10.1cm}@{}}
\toprule N.º & Observación & Estado & Cambio y alcance \\\midrule
\endfirsthead
\toprule N.º & Observación & Estado & Cambio y alcance \\\midrule
\endhead
1 & Proceso y retroalimentación & De proceso & No requiere edición; la retroalimentación externa no se certifica mediante archivos. \\[4pt]
2 & Resumen & Actualizado & Identifica corpus construido de 96 y ejecución de 42 pares; cifras actuales y alcance del avance. \\[4pt]
3 & Abstract & Actualizado & Se sincronizó con el resumen y los resultados de V3. \\[4pt]
4 & Estas/éstas & Revisado & Se normalizaron demostrativos del texto editado. No se dispone de la nota literal para acreditar el pasaje concreto. \\[4pt]
5 & Dos dimensiones & Atendido & Introducción: se distinguen formato y explicitud. \\[4pt]
6 & Definir conceptos desde el inicio & Atendido & Definiciones operativas de dirección, proyección, AUROC y formato en la introducción; formalización y glosario complementan. \\[4pt]
7 & Remisión a diseño & Atendido & Se atiende junto con 9 y 10. \\[4pt]
8 & Brecha y Pinto & Atendido en texto & Se conserva la distinción entre steering y diagnóstico; no se presenta como auditoría bibliográfica externa. \\[4pt]
9 & Cambio de instrucciones y herramientas & Corregido & Se separa variación entre pares de comparación pareada. La manipulación combina sistema y herramientas. \\[4pt]
10 & Simetría entre ramas & Verificado & 96/96 pares del corpus ampliado: mismo sistema, catálogo y esquemas dentro de cada formato. \\[4pt]
11 & Autoría y revisión asistida & Documentado & Se precisa composición por origen y responsabilidad del autor; revisión asistida no equivale a doble anotación. \\[4pt]
12 & Alcance de la proyección & Corregido & Se distinguen alineación geométrica, discriminación de contenido y rechazo conductual. \\[4pt]
13 & Dependencia y pareado & Corregido con límites & Interacción intra-par y evaluación retenida; se explicitan dependencia entre folds y supuestos entre pares. \\[4pt]
14 & Procedimiento y resultados PCA & Incorporado & Métodos y resultados: centrado global, PC1, varianza explicada y coseno; figura regenerada. \\[4pt]
15 & Extracción residual & Precisado & Último token del estímulo renderizado antes de generar; puede ser delimitador de plantilla. \\[4pt]
16 & Dirección y prueba de signos & Corregido & La prueba depende del eje; solo es invariante a reescalados positivos. Se distingue inferencia separada de cross-fitting. \\[4pt]
17 & Lineal/MLP y curva & Corregido & No se afirma equivalencia por cercanía. Se separan las métricas de sonda y dirección; las figuras usan la corrida actual. \\[4pt]
18 & 39/42 y selección de capa & Resuelto para 42 & 21 pares seleccionan y 21 infieren: Llama 21/21 decrecen, Mistral 19/21. La réplica de 96 permanece pendiente. \\[4pt]
19 & Conclusiones y recomendaciones & Redactado como avance & Se completaron discusión y conclusiones acotadas a los resultados disponibles; se habilitó el capítulo 6. \\[4pt]
20 & Material complementario & Incorporado & Anexos de corpus, configuración, evidencia, hashes y estado de observaciones. No reemplaza validación externa. \\[4pt]
\bottomrule\end{longtable}
\normalsize
\textbf{Pendientes experimentales.} Réplica sobre 96 pares, recuperación del paquete completo de Mistral V3, segunda anotación de explicitud, auditoría humana de generaciones y evaluación externa de monitores. Las conclusiones incorporadas son del avance, no un cierre de esas tareas.
\end{document}
